\chapter{Zbieżność Kuratowskiego oraz topologie Wijsmana i Attoucha--Wetsa}
\label{ch:topog22-zbieznosc-zbiorow}

\noindent\textbf{Cel rozdziału.}
W rozdziale~\ref{ch:topog-hiperprzestrzenie} wprowadziliśmy topologie
Vietorisa i Fella oraz metrykę Hausdorffa. Zbadamy teraz inne
sposoby opisywania granic zbiorów domkniętych: przez punkty,
które można przybliżać punktami kolejnych zbiorów, oraz przez funkcje
odległości od tych zbiorów.
Ustalimy dokładnie, kiedy te konstrukcje są równoważne.
Szczególną rolę odegra założenie, że domknięte kule są zwarte.
Nie zakładamy z góry, że rozważane zbiory domknięte są ograniczone.

\section{Granice dolna i górna Kuratowskiego}
\label{sec:topog22-kuratowski}

\begin{definicja}[Granice zbiorów w sensie Kuratowskiego]
\label{def:topog22-kuratowski}
Niech \((X,d)\) będzie przestrzenią metryczną
i niech \((F_n)_{n\ge1}\) będzie ciągiem jej podzbiorów.
Definiujemy \emph{granicę dolną Kuratowskiego}
oraz \emph{granicę górną Kuratowskiego} wzorami
\[
 \operatorname{Li}_{n\to\infty}F_n
 =\bigl\{x\in X:\ \forall\varepsilon>0\
     \exists N\ \forall n\ge N\
     F_n\cap B_d(x,\varepsilon)\neq\varnothing\bigr\},
\]
\[
 \operatorname{Ls}_{n\to\infty}F_n
 =\bigl\{x\in X:\ \forall\varepsilon>0\
     \forall N\ \exists n\ge N\
     F_n\cap B_d(x,\varepsilon)\neq\varnothing\bigr\}.
\]
Jeśli te zbiory są równe pewnemu \(F\subseteq X\),
mówimy, że \(F_n\) \emph{zbiega do \(F\) w sensie
Kuratowskiego}.
W tych definicjach dopuszczamy także zbiory puste.
\end{definicja}

\begin{stwierdzenie}[Opis przez funkcje odległości]
\label{prop:topog22-kuratowski-dist}
Niech \((X,d)\) będzie przestrzenią metryczną,
\((F_n)\) ciągiem podzbiorów \(X\), a
\(d(x,\varnothing)=+\infty\).
Wówczas
\[
 \operatorname{Li}F_n
 =\{x\in X:d(x,F_n)\longrightarrow0\},
 \qquad
 \operatorname{Ls}F_n
 =\{x\in X:\liminf_{n\to\infty}d(x,F_n)=0\}.
\]
Ponadto \(\operatorname{Li}F_n\subseteq\operatorname{Ls}F_n\),
a oba zbiory są domknięte w \(X\).
\end{stwierdzenie}
\begin{proof}
Dla dowolnego \(x\in X\) i \(\varepsilon>0\)
warunek \(F_n\cap B_d(x,\varepsilon)\neq\varnothing\)
jest równoważny \(d(x,F_n)<\varepsilon\).
Podstawiając tę równoważność do kwantyfikatorów
w Definicji~\ref{def:topog22-kuratowski}, otrzymujemy
oba wzory. Pierwszy warunek implikuje drugi.

Wykażemy domkniętość granicy dolnej bez założeń
o domkniętości \(F_n\).
Niech \(x\) należy do domknięcia \(\operatorname{Li}F_n\)
i niech \(\varepsilon>0\).
Wybierzmy \(y\in\operatorname{Li}F_n\)
z \(d(x,y)<\varepsilon/2\).
Od pewnego \(n\) zachodzi \(d(y,F_n)<\varepsilon/2\),
a z nierówności trójkąta
\(d(x,F_n)\le d(x,y)+d(y,F_n)<\varepsilon\).
Zatem \(x\in\operatorname{Li}F_n\).

Dla granicy górnej wybierzmy
\(x\in\overline{\operatorname{Ls}F_n}^{\,X}\).
Dla dowolnych \(\varepsilon>0\) i \(N\)
istnieje \(y\in\operatorname{Ls}F_n\)
z \(d(x,y)<\varepsilon/2\);
dla pewnego \(n\ge N\) mamy
\(d(y,F_n)<\varepsilon/2\).
Stąd \(d(x,F_n)<\varepsilon\).
Spełnione są wszystkie kwantyfikatory
definicji granicy górnej, więc
\(x\in\operatorname{Ls}F_n\).
\end{proof}

\begin{przyklad}[Brak granicy Kuratowskiego]
\label{ex:topog22-oscillating}
Niech \(X=\mathbb R\), \(F_{2n}=\{0\}\),
a \(F_{2n-1}=\{1\}\) dla \(n\ge1\).
Wówczas
\(\operatorname{Li}F_n=\varnothing\) oraz
\(\operatorname{Ls}F_n=\{0,1\}\).
Istotnie, żaden punkt nie ma wszystkich dostatecznie
późnych zbiorów w każdym swoim otoczeniu,
natomiast otoczenia \(0\) i \(1\) trafiają
nieskończenie wiele wyrazów ciągu.
Inne punkty nie należą do granicy górnej,
bo mają otoczenie rozłączne z \(\{0,1\}\).
\end{przyklad}

\begin{przyklad}[Punkt uciekający do nieskończoności]
\label{ex:topog22-escape}
W \(X=\mathbb R\) niech
\(F_n=\{0,n\}\).
Wtedy
\[
 \operatorname{Li}F_n=\operatorname{Ls}F_n=\{0\}.
\]
Dla \(x\neq0\) można bowiem wybrać kulę wokół \(x\)
rozłączną z \(\{0\}\), która od pewnego miejsca
nie zawiera również \(n\).
Nie zachodzi jednak zbieżność w metryce Hausdorffa,
ponieważ \(d_H(F_n,\{0\})=n\).
Przykład pokazuje potrzebę rozróżnienia
granicy lokalnej od globalnego pomiaru odległości zbiorów.

\begin{figure}[htbp]
\centering
\begin{tikzpicture}[x=1.08cm,y=.9cm,>=Latex]
 \draw[->,thin] (-.4,0)--(6.65,0);
 \foreach \x/\t in {0/0,1/1,2/2,3/3,4/4,5/5,6/6}
    {\draw (\x,-.08)--(\x,.08);
     \node[below,font=\scriptsize] at (\x,-.1) {$\t$};}
 \fill[blue!65!black] (0,.42) circle (2.5pt);
 \foreach \x in {2,4,6}
    {\fill[orange!80!black] (\x,.42) circle (2.5pt);}
 \node[above,blue!70!black,font=\small] at (0,.47) {punkt stały};
 \node[above,orange!85!black,font=\small] at (4,.47)
    {punkty oddalające się};
\end{tikzpicture}
\caption{Przy \(F_n=\{0,n\}\) punkt \(0\) pozostaje,
a drugi punkt opuszcza każdy ustalony zbiór ograniczony.}
\label{fig:topog22-escape}
\end{figure}
\end{przyklad}

\section{Topologia Wijsmana}
\label{sec:topog22-wijsman}

\begin{definicja}[Topologia Wijsmana]
\label{def:topog22-wijsman}
Niech \((X,d)\) będzie niepustą przestrzenią metryczną.
Przez \(\mathcal F_*(X)\) oznaczamy zbiór wszystkich
\emph{niepustych domkniętych} podzbiorów \(X\).
Dla \(x\in X\) definiujemy
\[
 \delta_x\colon\mathcal F_*(X)\to[0,\infty),
 \qquad \delta_x(F)=d(x,F).
\]
\emph{Topologią Wijsmana} na \(\mathcal F_*(X)\)
nazywamy najsłabszą topologię, względem której
wszystkie odwzorowania \(\delta_x\) są ciągłe.
W szczególności ciąg \(F_n\) jest w niej zbieżny
do \(F\in\mathcal F_*(X)\) dokładnie wtedy, gdy
\[
 d(x,F_n)\longrightarrow d(x,F)
 \quad\text{dla każdego }x\in X.
\]
\end{definicja}

\begin{uwaga}[Dlaczego wyłączamy zbiór pusty]
\label{rem:topog22-wijsman-empty}
Dla niepustego \(F\) funkcja \(d(\,\cdot\,,F)\)
ma wartości skończone, nawet gdy \(F\) jest nieograniczony.
Dla \(\varnothing\) przyjęliśmy wartość \(+\infty\).
Można osobno rozważać rozszerzoną topologię Wijsmana
z wartościami w \([0,\infty]\), lecz w tym rozdziale
porównania z metryką Attoucha--Wetsa dotyczą
wyłącznie \(\mathcal F_*(X)\).
Zbieżność w sensie Kuratowskiego i topologia Fella
nadal mogą obejmować \(\varnothing\).
\end{uwaga}

\begin{lemat}[Odległość od zbioru jest lipschitzowska]
\label{lem:topog22-distance-lipschitz}
Niech \((X,d)\) będzie przestrzenią metryczną
i \(F\subseteq X\) zbiorem niepustym.
Dla każdych \(x,y\in X\) zachodzi
\[
 |d(x,F)-d(y,F)|\le d(x,y).
\]
Jeśli \(F\) jest domknięty, to
\(F=\{x\in X:d(x,F)=0\}\).
\end{lemat}
\begin{proof}
Dla każdego \(z\in F\) mamy
\(d(x,F)\le d(x,z)\le d(x,y)+d(y,z)\).
Po wzięciu infimum względem \(z\in F\)
otrzymujemy \(d(x,F)\le d(x,y)+d(y,F)\).
Zamieniając \(x\) z \(y\), uzyskujemy pierwszą tezę.
Druga wynika z definicji domknięcia:
\(d(x,F)=0\) dokładnie wtedy, gdy
\(x\in\overline F^{\,X}\).
\end{proof}

\begin{twierdzenie}[Metryzowalność topologii Wijsmana]
\label{thm:topog22-wijsman-metrizable}
Niech \((X,d)\) będzie niepustą ośrodkową
przestrzenią metryczną i
\(D=\{q_1,q_2,\ldots\}\subseteq X\)
ustaloną przeliczalną rodziną gęstą
(dopuszczamy powtórzenia).
Wówczas wzór
\[
 \rho_W(A,B)=\sum_{j=1}^{\infty}2^{-j}
       \min\{1,\ |d(q_j,A)-d(q_j,B)|\},
 \qquad A,B\in\mathcal F_*(X),
\]
określa metrykę na \(\mathcal F_*(X)\),
której topologia jest topologią Wijsmana.
\end{twierdzenie}
\begin{proof}
Nieujemność i symetria są oczywiste,
a nierówność trójkąta wynika z
\(\min\{1,a+b\}\le\min\{1,a\}+\min\{1,b\}\)
dla \(a,b\ge0\).
Jeżeli \(\rho_W(A,B)=0\), to
\(d(q_j,A)=d(q_j,B)\) dla każdego \(j\).
Z Lematu~\ref{lem:topog22-distance-lipschitz}
obie funkcje odległości są ciągłe,
a \(D\) jest gęsty, zatem są równe na całym \(X\).
Ponieważ \(A,B\) są domknięte, ich zbiory zerowe
są równe, więc \(A=B\).

Zbieżność w metryce \(\rho_W\) jest równoważna
zbieżności \(d(q_j,F_n)\to d(q_j,F)\) dla każdego \(j\):
w jedną stronę wynika to z oszacowania
pojedynczego składnika, a w drugą ze zbieżności
ograniczonego szeregu, rozbijanego na
skończenie wiele pierwszych składników i mały ogon.
Dla dowolnego \(x\in X\), \(q_j\in D\)
i \(A,B\in\mathcal F_*(X)\) mamy
\[
 |d(x,A)-d(x,B)|
 \le 2d(x,q_j)+|d(q_j,A)-d(q_j,B)|.
\]
Dzięki gęstości \(D\) zbieżność na \(D\)
pociąga zbieżność na całym \(X\).
Dowodzi to równoważności z topologią Wijsmana.
Ten sam argument z otoczeniami zamiast ciągów
pokazuje równość topologii, a nie tylko
ich zbieżnych ciągów.
\end{proof}

\begin{twierdzenie}[Wijsman implikuje Kuratowskiego]
\label{thm:topog22-wijsman-kuratowski}
Niech \((X,d)\) będzie przestrzenią metryczną
i \(F,F_n\in\mathcal F_*(X)\) dla \(n\ge1\).
Jeżeli \(F_n\to F\) w topologii Wijsmana,
to \(F_n\to F\) w sensie Kuratowskiego.
\end{twierdzenie}
\begin{proof}
Dla każdego \(x\in F\) mamy \(d(x,F)=0\),
a więc \(d(x,F_n)\to0\).
Stwierdzenie~\ref{prop:topog22-kuratowski-dist}
daje \(x\in\operatorname{Li}F_n\).
Jeżeli \(x\notin F\), domkniętość \(F\) zapewnia
\(d(x,F)>0\). Stąd dla dostatecznie dużych \(n\)
zachodzi \(d(x,F_n)>d(x,F)/2\), więc
\(x\notin\operatorname{Ls}F_n\).
Otrzymujemy
\(F\subseteq\operatorname{Li}F_n
\subseteq\operatorname{Ls}F_n\subseteq F\).
\end{proof}

\section{Topologia Attoucha--Wetsa}
\label{sec:topog22-aw}

\begin{definicja}[Metryka i topologia Attoucha--Wetsa]
\label{def:topog22-aw}
Niech \((X,d)\) będzie niepustą przestrzenią
metryczną i niech \(x_0\in X\).
Dla \(m\ge1\) oraz \(A,B\in\mathcal F_*(X)\) połóżmy
\[
 \Delta_m(A,B)=
  \sup_{x\in\overline B_d(x_0,m)}
       |d(x,A)-d(x,B)|,
\]
\[
 \rho_{AW}(A,B)=\sum_{m=1}^{\infty}
          2^{-m}\min\{1,\Delta_m(A,B)\}.
\]
Topologię tej metryki nazywamy
\emph{topologią Attoucha--Wetsa}.
Zbiory \(\overline B_d(x_0,m)\)
są jedynie kulami domkniętymi:
w definicji nie zakładamy ich zwartości.
\end{definicja}

\begin{twierdzenie}[Metryka Attoucha--Wetsa i zbieżność na ograniczonych zbiorach]
\label{thm:topog22-aw-metric}
Niech \((X,d)\) będzie niepustą przestrzenią metryczną,
\(x_0\in X\), a \(\mathcal F_*(X)\) rodziną
niepustych domkniętych podzbiorów.
Funkcja \(\rho_{AW}\) z Definicji~\ref{def:topog22-aw}
jest metryką. Ciąg \((F_n)\subseteq\mathcal F_*(X)\)
zbiega w niej do \(F\in\mathcal F_*(X)\)
dokładnie wtedy, gdy dla każdego ograniczonego
zbioru \(B\subseteq X\) zachodzi
\[
 \sup_{x\in B}|d(x,F_n)-d(x,F)|\longrightarrow0.
\]
Topologia ta nie zależy od wyboru punktu \(x_0\).
\end{twierdzenie}
\begin{proof}
Dla ustalonych \(A,B\in\mathcal F_*(X)\)
oraz \(x\in\overline B_d(x_0,m)\) mamy
\(d(x,A)\le m+d(x_0,A)\) i analogicznie dla \(B\).
Zatem \(\Delta_m(A,B)\) jest skończona.
Nieujemność i symetria są oczywiste;
nierówność trójkąta dla \(\Delta_m\) wynika
z nierówności trójkąta dla wartości bezwzględnej,
a dla \(\rho_{AW}\) z nierówności
\(\min\{1,a+b\}\le\min\{1,a\}+\min\{1,b\}\).

Jeżeli \(\rho_{AW}(A,B)=0\), to wszystkie
\(\Delta_m(A,B)\) są zerowe.
Ponieważ każde \(x\in X\) należy do pewnej
\(\overline B_d(x_0,m)\), funkcje odległości
do \(A\) i do \(B\) są identyczne.
Lemat~\ref{lem:topog22-distance-lipschitz}
daje \(A=B\).

Z definicji szeregu wynika, że
\(\rho_{AW}(F_n,F)\to0\) wtedy i tylko wtedy,
gdy \(\Delta_m(F_n,F)\to0\) dla każdego \(m\).
Każdy ograniczony zbiór \(B\) zawiera się
w pewnej kuli \(\overline B_d(x_0,m)\),
więc otrzymujemy żądaną zbieżność jednostajną.
Implikacja odwrotna wynika bezpośrednio
z zastosowania założenia do każdej takiej kuli.
Ponieważ rodziny kul wokół dwóch różnych
punktów \(x_0,x_1\) są współkońcowe
w rodzinie wszystkich ograniczonych podzbiorów \(X\),
otrzymana topologia nie zależy od punktu bazowego.
\end{proof}

\begin{stwierdzenie}[Attouch--Wets jest silniejsza od Wijsmana]
\label{prop:topog22-aw-wijsman}
Dla każdej niepustej przestrzeni metrycznej \((X,d)\)
topologia Attoucha--Wetsa na \(\mathcal F_*(X)\)
jest co najmniej tak silna jak topologia Wijsmana.
\end{stwierdzenie}
\begin{proof}
Każdy singleton \(\{x\}\subseteq X\) jest ograniczony.
Zbieżność jednostajna funkcji odległości
na wszystkich ograniczonych zbiorach
pociąga zbieżność w każdym punkcie.
Ponadto każda ewaluacja
\(F\mapsto d(x,F)\) jest ciągła względem
topologii Attoucha--Wetsa, ponieważ
\[
 |d(x,A)-d(x,B)|\le\Delta_m(A,B)
\]
dla \(m\) takiego, że \(x\in\overline B_d(x_0,m)\).
Z najsłabszości topologii Wijsmana
wynika wymagane zawieranie topologii.
\end{proof}

\begin{przyklad}[Ścisłość na ograniczonej przestrzeni]
\label{ex:topog22-strict-aw}
Niech \(X=\mathbb N\) z metryką dyskretną:
\(d(k,l)=1\) dla \(k\neq l\).
Jest to przestrzeń metryczna ograniczona,
zupełna i ośrodkowa, ale niezwarta.
Niech \(F_n=X\setminus\{n\}\) i \(F=X\).
Dla każdego ustalonego \(k\in X\)
od pewnego \(n\) zachodzi \(k\in F_n\);
wtedy \(d(k,F_n)=0=d(k,F)\).
Zatem \(F_n\to F\) w topologii Wijsmana.

Jednak \(d(n,F_n)=1\), podczas gdy
\(d(n,F)=0\). Stąd
\[
 \sup_{x\in X}|d(x,F_n)-d(x,F)|=1
 \quad\text{dla każdego }n.
\]
Ponieważ \(X\) jest ograniczony,
nie ma zbieżności w topologii Attoucha--Wetsa.
Obie topologie są więc różne nawet na rodzinie
niepustych zbiorów domkniętych zupełnej,
ośrodkowej przestrzeni metrycznej.
\end{przyklad}

\section{Topologia Fella i właściwe przestrzenie metryczne}
\label{sec:topog22-proper}

\begin{definicja}[Przestrzeń metryczna właściwa]
\label{def:topog22-proper}
Przestrzeń metryczną \((X,d)\) nazywamy
\emph{właściwą}, jeżeli każda domknięta
kula \(\overline B_d(x,r)\), gdzie
\(x\in X\) i \(r\ge0\), jest zwarta.
Wówczas każdy domknięty ograniczony
podzbiór \(X\) jest zwarty.
\end{definicja}

\begin{lemat}[Zbieżność funkcji odległości na zwartym zbiorze]
\label{lem:topog22-equilipschitz}
Niech \((X,d)\) będzie przestrzenią metryczną,
\(K\subseteq X\) zbiorem zwartym,
a \((f_n)\) ciągiem funkcji
\(f_n\colon X\to\mathbb R\)
spełniających
\(|f_n(x)-f_n(y)|\le d(x,y)\) dla wszystkich \(x,y\in X\).
Jeżeli \(f_n(x)\to f(x)\) dla każdego \(x\in X\),
to \(f_n\to f\) jednostajnie na \(K\).
\end{lemat}
\begin{proof}
Przechodząc do granicy w nierówności,
stwierdzamy, że \(f\) również jest
\(1\)-lipschitzowska.
Jeśli \(K=\varnothing\), nie ma czego dowodzić.
Niech \(\varepsilon>0\).
Zwartość \(K\) daje punkty \(x_1,\ldots,x_m\in K\)
takie, że każdy \(x\in K\)
spełnia \(d(x,x_i)<\varepsilon/3\)
dla pewnego \(i\).
Dla dostatecznie dużych \(n\) mamy
\(|f_n(x_i)-f(x_i)|<\varepsilon/3\)
jednocześnie dla wszystkich \(i\).
Wtedy
\[
 |f_n(x)-f(x)|
 \le |f_n(x)-f_n(x_i)|
     +|f_n(x_i)-f(x_i)|+|f(x_i)-f(x)|
 <\varepsilon.
\]
\end{proof}

\begin{twierdzenie}[Równoważność Wijsmana i Attoucha--Wetsa]
\label{thm:topog22-proper-wijsman-aw}
Niech \((X,d)\) będzie niepustą właściwą
przestrzenią metryczną. Wówczas
topologia Wijsmana na \(\mathcal F_*(X)\)
pokrywa się z topologią Attoucha--Wetsa.
\end{twierdzenie}
\begin{proof}
Zawieranie topologii Wijsmana w Attoucha--Wetsa
wynika ze Stwierdzenia~\ref{prop:topog22-aw-wijsman}.
Odwrotnie, dla dowolnej zbieżnej w sensie Wijsmana
sieci \((F_\lambda)\) do \(F\)
funkcje \(d(\,\cdot\,,F_\lambda)\)
są \(1\)-lipschitzowskie i zbieżne punktowo
do \(d(\,\cdot\,,F)\).
Każda kula \(\overline B_d(x_0,m)\) jest zwarta.
Argument skończonej \(\varepsilon\)-sieci
z Lematu~\ref{lem:topog22-equilipschitz}
działa również dla sieci:
dla skończenie wielu punktów bazowych
wybiera się wspólny indeks dominujący.
Zatem zbieżność jest jednostajna na każdej
kuli i na każdym zbiorze ograniczonym.
Twierdzenie~\ref{thm:topog22-aw-metric}
pokazuje, że jest to zbieżność
w topologii Attoucha--Wetsa.
W szczególności identyczność z topologii Wijsmana
do topologii Attoucha--Wetsa jest ciągła.
\end{proof}

\begin{twierdzenie}[Zbieżność Fella i Kuratowskiego]
\label{thm:topog22-fell-kuratowski}
Niech \((X,d)\) będzie lokalnie zwartą
przestrzenią metryczną, a \(F,F_n\subseteq X\)
zbiorami domkniętymi, przy czym dopuszczamy
\(F=\varnothing\) i \(F_n=\varnothing\).
Wówczas \(F_n\to F\) w topologii Fella
z Definicji~\ref{def:topog21-fell}
dokładnie wtedy, gdy
\[
 \operatorname{Li}F_n=\operatorname{Ls}F_n=F.
\]
\end{twierdzenie}
\begin{proof}
Niech najpierw \(F_n\to F\) w topologii Fella.
Dla każdego \(x\in F\) i \(\varepsilon>0\)
otwarty zbiór \(B_d(x,\varepsilon)\) trafia \(F\),
więc na mocy warunku dolnego
zbieżności Fella jest od pewnego miejsca
trafiany przez wszystkie \(F_n\).
Stąd \(F\subseteq\operatorname{Li}F_n\).

Niech \(x\notin F\).
Z lokalnej zwartości i własności Hausdorffa
oraz Lematu~\ref{lem:topog-baire-zwarte-domkniecie}
istnieje otwarte \(W\ni x\)
o zwartym domknięciu \(C\subseteq X\setminus F\).
Warunek ominięcia zwartego \(C\)
zapewnia \(F_n\cap C=\varnothing\)
od pewnego miejsca, a zatem
\(F_n\cap W=\varnothing\) od pewnego miejsca.
Stąd \(x\notin\operatorname{Ls}F_n\).
Otrzymujemy
\(F\subseteq\operatorname{Li}F_n
\subseteq\operatorname{Ls}F_n\subseteq F\).

Odwrotnie, niech obie granice Kuratowskiego
są równe \(F\).
Jeżeli otwarty \(U\) trafia \(F\),
wybierzmy \(x\in F\cap U\) i
\(\varepsilon>0\) tak, aby
\(B_d(x,\varepsilon)\subseteq U\).
Z \(x\in\operatorname{Li}F_n\)
wynika, że \(F_n\) ostatecznie trafiają \(U\).

Niech \(C\subseteq X\) będzie zwarty
i rozłączny z \(F\).
Gdyby nieskończenie wiele \(F_n\) trafiało \(C\),
moglibyśmy wybrać rosnący ciąg indeksów \(n_k\)
oraz \(x_k\in F_{n_k}\cap C\).
Zwartość metryczna \(C\) daje podciąg
\(x_{k_j}\to x\in C\).
Dla każdego \(\varepsilon>0\) i \(N\)
istnieje \(j\) z \(n_{k_j}\ge N\)
oraz \(d(x_{k_j},x)<\varepsilon\).
To oznacza \(x\in\operatorname{Ls}F_n=F\),
wbrew \(C\cap F=\varnothing\).
Zatem \(F_n\) ostatecznie omijają \(C\).
Oba warunki zbieżności Fella są spełnione
na mocy Stwierdzenia~\ref{prop:topog21-fell-nets}.
\end{proof}

\begin{twierdzenie}[Zbieżności są równoważne w przestrzeni właściwej]
\label{thm:topog22-proper-equivalences}
Niech \((X,d)\) będzie niepustą właściwą
przestrzenią metryczną oraz
\(F,F_n\in\mathcal F_*(X)\) dla wszystkich \(n\ge1\).
Następujące warunki są równoważne:
\begin{enumerate}[label=\textup{(\roman*)}]
\item \(F_n\to F\) w sensie Kuratowskiego;
\item \(F_n\to F\) w topologii Fella;
\item \(F_n\to F\) w topologii Wijsmana;
\item \(F_n\to F\) w topologii Attoucha--Wetsa.
\end{enumerate}
\end{twierdzenie}
\begin{proof}
Przestrzeń właściwa jest lokalnie zwarta,
ponieważ kula domknięta o dodatnim promieniu
jest zwartym otoczeniem swojego środka.
Twierdzenie~\ref{thm:topog22-fell-kuratowski}
daje \textup{(i)} \(\Leftrightarrow\) \textup{(ii)}.
Twierdzenia~\ref{thm:topog22-wijsman-kuratowski}
i~\ref{thm:topog22-proper-wijsman-aw}
oraz Stwierdzenie~\ref{prop:topog22-aw-wijsman}
dają implikacje \textup{(iv)} \(\Rightarrow\)
\textup{(iii)} \(\Rightarrow\) \textup{(i)}
oraz \textup{(iii)} \(\Leftrightarrow\) \textup{(iv)}.
Pozostaje dowieść \textup{(i)} \(\Rightarrow\) \textup{(iii)}.

Ustalmy \(x\in X\) i połóżmy \(r=d(x,F)\).
Ponieważ \(F\neq\varnothing\), liczba \(r\) jest skończona.
Dla każdego \(\varepsilon>0\) wybierzmy \(z\in F\)
takie, że \(d(x,z)<r+\varepsilon\).
Z \(z\in\operatorname{Li}F_n\)
wynika \(d(z,F_n)\to0\).
Stąd
\[
 \limsup_{n\to\infty}d(x,F_n)
 \le d(x,z)<r+\varepsilon.
\]
Po przejściu z \(\varepsilon\downarrow0\)
dostajemy \(\limsup_nd(x,F_n)\le r\).

Jeżeli \(r=0\), nieujemność odległości
kończy dowód. Niech \(r>0\).
Gdyby \(\liminf_nd(x,F_n)<r\),
istniałyby \(\eta>0\) i rosnące indeksy \(n_k\)
takie, że
\(d(x,F_{n_k})<r-\eta\).
Wybierzmy \(z_k\in F_{n_k}\)
z \(d(x,z_k)<r-\eta/2\).
Wszystkie \(z_k\) należą do zwartej
kuli \(\overline B_d(x,r)\);
po przejściu do podciągu możemy
przyjąć \(z_k\to z\).
Ponieważ \(n_k\to\infty\)
i \(z_k\in F_{n_k}\),
punkt \(z\) należy do
\(\operatorname{Ls}F_n=F\).
Jednocześnie \(d(x,z)\le r-\eta/2<r\),
wbrew definicji \(r=d(x,F)\).
Stąd \(\liminf_nd(x,F_n)\ge r\)
i ostatecznie \(d(x,F_n)\to d(x,F)\).
Punkt \(x\) był dowolny,
więc otrzymujemy topologię Wijsmana.
\end{proof}

\begin{przyklad}[Lokalna zwartość nie wystarcza dla Wijsmana]
\label{ex:topog22-nonproper}
Niech \(X=\{p,q\}\cup\{y_n:n\ge1\}\)
i niech odległości różnych punktów spełniają
\[
 d(p,q)=2,\qquad
 d(p,y_n)=d(q,y_n)=1,\qquad
 d(y_n,y_m)=1\quad(n\neq m).
\]
Są to odległości metryczne:
jedyny odcinek długości \(2\)
ma końce \(p,q\), a każda droga
łącząca je przez trzeci punkt ma długość \(2\).
Przestrzeń jest dyskretna, przeliczalna,
ograniczona, zupełna i lokalnie zwarta,
lecz nie jest właściwa, ponieważ sama
jest domkniętą kulą promienia \(2\)
i nie jest zwarta.

Niech \(F_n=\{p,y_n\}\), \(F=\{p\}\).
Każdy punkt \(y_k\) należy tylko do jednego \(F_n\),
a \(p\) należy do wszystkich.
Ponieważ topologia \(X\) jest dyskretna,
\(\operatorname{Li}F_n=\operatorname{Ls}F_n=\{p\}\).
Stąd \(F_n\to F\) również w topologii Fella
na mocy Twierdzenia~\ref{thm:topog22-fell-kuratowski}.
Jednak
\[
 d(q,F_n)=1,\qquad d(q,F)=2
 \quad\text{dla każdego }n,
\]
więc nie ma zbieżności Wijsmana.
Pokazuje to, dlaczego w
Twierdzeniu~\ref{thm:topog22-proper-equivalences}
potrzebna jest dodatkowa zwartość kul,
a sama lokalna zwartość nie wystarcza.
\end{przyklad}

\begin{przyklad}[Zbieżność lokalna nie oznacza zbieżności Hausdorffa]
\label{ex:topog22-aw-not-hausdorff}
W przestrzeni właściwej \(X=\mathbb R\)
niech \(F_n=\{0,n\}\) i \(F=\{0\}\).
Dla każdej ustalonej kuli \([-m,m]\)
oraz \(n>2m\) mamy
\(d(x,F_n)=|x|=d(x,F)\) dla \(x\in[-m,m]\).
Zatem \(\Delta_m(F_n,F)=0\) od pewnego miejsca
i \(F_n\to F\) w topologii Attoucha--Wetsa.
Jednocześnie \(d_H(F_n,F)=n\).
W konsekwencji nawet na rodzinie
niepustych zbiorów zwartych w \(\mathbb R\)
topologia Vietorisa z rozdziału~21
jest silniejsza od topologii Attoucha--Wetsa.
\end{przyklad}

\section{Podsumowanie i zakres twierdzeń}
\label{sec:topog22-summary}

Granice Kuratowskiego wymagają jedynie struktury metrycznej
i opisują punkty dające się przybliżać punktami kolejnych zbiorów.
Topologia Wijsmana bada punktową zbieżność funkcji odległości
do niepustych zbiorów domkniętych,
a topologia Attoucha--Wetsa --- ich zbieżność jednostajną
na zbiorach ograniczonych.
Ta ostatnia jest zawsze silniejsza od Wijsmana;
właściwość przestrzeni metrycznej czyni je równymi.
Dla lokalnie zwartych przestrzeni metrycznych
zbieżność ciągów w topologii Fella jest równoważna
zbieżności w sensie Kuratowskiego.
W przestrzeniach właściwych wszystkie cztery
kryteria są równoważne dla ciągów zbiorów
niepustych domkniętych o niepustej granicy.

Nie są to jednak synonimy metryki Hausdorffa:
punkt, który opuszcza każdy zbiór ograniczony,
może zanikać w granicy Kuratowskiego, Fella,
Wijsmana i Attoucha--Wetsa, chociaż odległość
Hausdorffa do granicy rośnie bez ograniczenia.
Zbiory puste dopuszczamy w granicach Kuratowskiego
i w topologii Fella, lecz nie w przyjętych
tu definicjach Wijsmana i Attoucha--Wetsa.